% Preserved fragment; see MANIFIESTO_ANTECEDENTES_LECTURA.json.
\section{Coinductive generation of \texorpdfstring{\(\pi,\varphi,e\)}{pi, phi, e} to arbitrary depth}
\label{sec:generacion}

The six-block emission is a finite stage of the construction, not a complete decimal expansion. Its role is to produce regions with multiplicity, orientation, and transport rules. Prolongation must preserve these data and determine a compatible collection of prefixes of arbitrary length. The term \emph{generation} denotes this operation; \emph{genealogy} denotes the reconstruction of its dependencies. Both notions are involved, but they serve different purposes.

\subsection{Orbits, regions, and characters of closure, propagation, and self-scaling}

For a word \(w=(w_1,\ldots,w_6)\), define
\[
 r(w)=(w_3,w_1,w_2,w_5,w_6,w_4),\qquad
 s(w)=(w_4,w_5,w_6,w_1,w_2,w_3).
\]
We have \(r^3=s^2=I\) and \(srs=r^{-1}\); these permutations therefore realize the dihedral group \(D_3\). The action preserves the catalogue and the multiplicities and commutes with \(\Psi\). The quotient of the \(243\) visible words has \(43\) orbits: thirty-eight of cardinality six and five of cardinality three.

The closure class is recognized by a fibre predicate: its six orientations have five emissions per word, total multiplicity \(1008\), and distribution \(432+4\cdot144\). There is exactly one orbit with these properties in the catalogue. To specify the other two selections, write \(f(w)=\#\Psi^{-1}(w)\), \(m(w)=\sum_{U\in\Psi^{-1}(w)}\operatorname{mult}(U)\), and \(\nu(w)=(n_0,n_1,n_2)\). The additive diagonal class of an emission is \(d(U)=[i+j]_9\) in indices from zero to eight; its additive signature retains this class and the orientation \(ES\) or \(NO\). Let \(\mathcal D(\mathcal O)\) be the set of these classes for all emissions over an orbit.

Among the orbits of size six, the joint predicate
\[
 f(w)=1,\qquad m(w)=144\quad(w\in\mathcal O),\qquad
 \mathcal D(\mathcal O)=\{0,3,6\}
\]
leaves exactly six orbits. In this sector, the census \(\nu=(2,3,1)\), the same as for closure, uniquely selects propagation; the census \(\nu=(2,2,2)\) uniquely selects self-scaling. The diagonal support cannot be omitted: without it there are four balanced single-preimage orbits of multiplicity \(144\), not one. Finally, the existence of an emission with \(d=6\) and orientation \(ES\) selects one representative of each orbit and yields
\begin{equation}
 w^{\rm cl}=010211,\qquad w^{\rm pr}=201101,\qquad w^{\rm au}=121200.
 \label{eq:palabras-regionales}
\end{equation}
These symbols first designate regions of the catalogue. Scalar identification will be established through their readouts; they are not chosen because they coincide with digits of a known constant.

The closure word \(010211\) has the following five preimages:
\begin{center}\small
\begin{tabular}{@{}lll r@{}}\toprule
Six-block emission&Multiplicative signature&Role&Multiplicity\\\midrule
\(501,614,498,169,272,272\)&\(555555\)&transverse&432\\
\(810,923,870,169,272,272\)&\(339555\)&positive orientation&144\\
\(810,983,810,169,272,272\)&\(393555\)&positive orientation&144\\
\(870,923,810,169,272,272\)&\(933555\)&positive orientation&144\\
\(870,923,810,418,674,521\)&\(933717\)&oriented return&144\\\bottomrule
\end{tabular}
\end{center}
The multiplicity \(432\) and the constant signature \(555555\) characterize the transverse region. They do not identify the central position \((5,5)\) of APP. Confusing a six-window signature with a position in the support would erase the distinction between numerical region and electronic structure.

The microfibre also has a precise bipartite morphology. Each emission is a product of an additive cursor class and a multiplicative class. The three positive regions share the same additive class of eighteen conditions; together they contain three multiplicative classes of eight conditions each. The transverse region has a product \(18\times24\); the three positive regions taken together, another \(18\times24\); and the return, a product \(18\times8\). The decomposition by emissions has five parts, whereas the decomposition into connected components of the bipartite graph has three. Both preserve the same \(1008\) seeds without reducing form to a cardinality.

\subsection{Finite transport and the prolongation boundary}

The endomorphisms \(L_0,L_1\) of the ternary word space act before the Archimedean realization. Their determination begins with the transitions of the arithmetic state, not with the prescription of two matrices. If \(U_t\) is the composite step of selection, transport, and update, the observed block satisfies
\[
 b_j^\chi=\Pi_6(x_j^\chi),\qquad
 x_{j+1}^\chi=U_{9j+9}\cdots U_{9j+1}(x_j^\chi),
 \qquad \chi\in\{\mathrm{cl},\mathrm{pr},\mathrm{au}\}.
\]
Six transitions are formed for each regime: two consecutive transitions in each of the three regions. The pairs of source and target matrices obtained in the record are
\[
\begin{array}{c|c|c|c}
X_0&Y_0&X_1&Y_1\\\hline
010211&012222&010211&002111\\
012222&010211&002111&110221\\
201101&121221&102011&012222\\
121221&102011&012222&102011\\
121200&112202&121020&010210\\
112202&121020&010210&010200
\end{array}
\]
Each string of six digits represents a row in \(\FF_3^6\). Since \(\det X_0=1\) and \(\det X_1=-1\) in \(\FF_3\), the equations \(X_aL_a=Y_a\) uniquely determine \(L_a=X_a^{-1}Y_a\). In this row-vector chart, they yield
\[
 L_0=\begin{pmatrix}
 2&2&2&1&2&1\\2&1&2&2&1&1\\1&0&0&1&0&2\\
 0&0&1&1&1&0\\2&2&2&0&2&1\\2&1&2&1&0&0
 \end{pmatrix},\quad
 L_1=\begin{pmatrix}
 2&1&2&0&2&2\\1&2&1&1&2&0\\1&2&1&1&1&2\\
 0&1&0&0&2&1\\2&0&2&1&0&1\\0&2&2&2&1&1
 \end{pmatrix}\quad(\bmod3).
\]
The schedule \((L_0,L_0,L_1,L_1)\) produces four new blocks of six components. Their concatenations form
\[
 w_6\longrightarrow w_{12}\longrightarrow w_{18}
 \longrightarrow w_{24}\longrightarrow w_{30}.
\]
The subscripts of \(w\) indicate accumulated length. The boundary \(R_{36}\) preserves the terminal structure and opens the prolongation relation; it is not renamed as a final periodic block that could be repeated indefinitely.

The identity \(L=X^{-1}Y\) proves uniqueness once the transitions are fixed; it does not replace the production of those transitions by \(U_t\). The provenance of their record and the examination of the sixteen binary schedules are documented in the technical supplement. The indicated schedule is the only one of these sixteen that jointly reproduces the three sequences of five blocks.

\subsection{Five closure regions and angular compensation}

The closure region contains one transverse component and four accompanying components. The symmetric readout acts on the space of its five positions through
\[
 P_5=\frac15\mathbf1\mathbf1^{\mathsf T},\qquad
 \mathbf1=(1,1,1,1,1)^{\mathsf T}.
\]
Invariance \(P_5\lambda=\lambda\) and conservation \(\mathbf1^{\mathsf T}\lambda=1\) force \(\lambda=\mathbf1/5\). This normalization distributes the unit among positions of the readout; it does not replace the seed multiplicities \(432,144,144,144,144\) with uniform frequencies. The elliptic chart of the readout realizes each coordinate through \(1+\lambda_jJ\); its slope is \(\lambda_j\), so the primitive slope is \(1/5\). This makes explicit the rule taking a normalized coordinate to a slope, rather than identifying the two notions without a map. The readout composes the four accompanying positions by the law
\begin{equation}
 a\oplus b=\frac{a+b}{1-ab}.
 \label{eq:suma-pendientes}
\end{equation}
This law is obtained by multiplying \((1+aJ)(1+bJ)\) in the regime \(J^2=-I\) and dividing by its scalar component. The composition of slopes is therefore a realization of the quadratic structure of TRIT, not an ordinary sum of four fractions.

\begin{proposition}[Closure compensator]
The composition of four slopes \(1/5\) is \(120/119\). The condition of return to slope one determines a unique positive compensator of the form \(1/q\), with \(q>1\), and yields \(q=239\).
\end{proposition}
\begin{proof}
The first doubling gives \((1/5)\oplus(1/5)=5/12\); the second, \((5/12)\oplus(5/12)=120/119\). The condition
\[
 \frac{120/119-1/q}{1+(120/119)/q}=1
\]
is equivalent to \((120-119)q=120+119\), hence \(q=239\).
\end{proof}

The resulting closure character has the realization
\begin{equation}
 \Pi_{\rm cl}=4\left(4\arctan\frac15-\arctan\frac1{239}\right).
 \label{eq:lector-pi}
\end{equation}
Machin's angular identity recognizes \(\Pi_{\rm cl}=\pi\). The direction of the construction is precise: the slope and the number of compositions are specified by the pentafibre readout; the compensation equation forces \(239\). Once this character has been produced, its evaluation by alternating series does not select the seed orbit anew.

For \(0<x<1\), the alternating sums
\[
 A_n(x)=\sum_{j=0}^n\frac{(-1)^jx^{2j+1}}{2j+1}
\]
enclose \(\arctan x\) between two consecutive truncations, with difference \(x^{2n+3}/(2n+3)\). Applied to the two slopes in \eqref{eq:lector-pi}, they provide rational endpoints with explicit error bounds. The numerical value is thus calculated as a consequence of the exact character, without introducing a target prefix.

\subsection{Propagation readout and normalization of the unit}

The propagation region is realized by a formal collection \(P_t\in\mathcal A[[t]]\), where \(\mathcal A\) is a unital commutative algebra over \(\QQ\). Its composition satisfies \(P_{t+s}=P_tP_s\), with unit \(P_0=1\), and the oriented elementary transport fixes the scalar generator \(P'_0=+1\). The condition fixes its sign, not merely its magnitude. Writing
\[
 P_t=\sum_{n\ge0}a_nt^n,
\]
the normalization determines \(a_0=a_1=1\). Comparing the coefficient of \(s^nt\) in \(P_{s+t}=P_sP_t\) gives \((n+1)a_{n+1}=a_na_1\); equivalently, \(P'_t=P_t\). Thus,
\begin{equation}
 (n+1)a_{n+1}=a_n,\qquad a_n=\frac1{n!}.
 \label{eq:propagacion}
\end{equation}
The character at one unit of the parameter is \(\Pi_{\rm pr}=P_1\), subsequently recognized as \(e\).

Normalization of the generator matters. Without this datum, the semigroup law would admit \(P_t=e^{ct}\) for different \(c\); the condition fixing the unit of the parameter in the readout is therefore not omitted. This condition precedes decimal comparison and belongs to the exposition of the operational rule, not to the choice of an experimental value.

\begin{proposition}[Rational propagation bounds]
For \(n\ge1\), let \(E_n=\sum_{j=0}^n1/j!\). Then
\[
 E_n<\Pi_{\rm pr}<E_n+\frac{n+2}{(n+1)(n+1)!}.
\]
\end{proposition}
\begin{proof}
The first omitted term is \(1/(n+1)!\). Each subsequent ratio is at most \(1/(n+2)\). The geometric upper bound for the tail is
\[
 \frac1{(n+1)!}\frac1{1-1/(n+2)}
 =\frac{n+2}{(n+1)(n+1)!}.
\]
Strict inequality follows from the decrease of the successive ratios.
\end{proof}

\subsection{Modal incidence and generation of self-scaling}

The tritic selector produces the cycle \((+1,-1,0)\). On the modes \(\Sigma,\Pi\), the rule is \(+1:m\mapsto\Sigma\), \(-1:m\mapsto\Pi\), \(0:m\mapsto m\). With cyclic closure, the observed mode is \((\Sigma,\Pi,\Pi)\). Its edges are \(\Sigma\to\Pi\), \(\Pi\to\Pi\), and \(\Pi\to\Sigma\). Its incidence, in this ordering of the modes, is therefore
\begin{equation}
 F_{\rm av}=\begin{pmatrix}0&1\\1&1\end{pmatrix}.
 \label{eq:fav}
\end{equation}
The four entries are obtained from the schedule, not from an equation that has already received the value of \(\varphi\). Repeating the schedule gives counts \(3F_{\rm av},9F_{\rm av},36F_{\rm av}\) at lengths nine, twenty-seven, and one hundred and eight. These counts are not powers of the matrix.

\begin{theorem}[Positive self-scaling character]
The incidence \eqref{eq:fav} has a unique strictly positive eigenray. When normalized as \((1,x)^{\mathsf T}\), its coordinate \(x\) satisfies \(x^2-x-1=0\), \(x>0\), and equals \(\varphi=(1+\sqrt5)/2\).
\end{theorem}
\begin{proof}
The eigenvalue equation gives \(x=\lambda\) and \(1+x=\lambda x\). Eliminating \(\lambda\) yields the polynomial. Only one of its two roots is positive. Since all entries of \(F_{\rm av}^2\) are positive, the positive ray is unique and simple. The radical expression recognizes the coordinate obtained; it does not participate in the production of \(F_{\rm av}\).
\end{proof}

With \(F_0=0,F_1=1,F_{n+1}=F_n+F_{n-1}\), the incidence induces, for \(n\geq1\),
\[
 F_{\rm av}^n=\begin{pmatrix}F_{n-1}&F_n\\F_n&F_{n+1}\end{pmatrix}.
\]
The consecutive ratios \(F_{n+1}/F_n\) alternate around \(\varphi\). The determinant identity
\[
 F_{n+1}^2-F_nF_{n+2}=(-1)^n
\]
gives an exact width \(1/(F_nF_{n+1})\) between two consecutive ratios. Their growth makes this rational bracket arbitrarily narrow.

\subsection{Compatible cylinders and arbitrary depth}

The realization of a ternary block \(b=(b_1,\ldots,b_6)\) is the integer
\[
 \nu(b)=\sum_{j=1}^6b_j3^{6-j}\in\{0,\ldots,728\}.
\]
A sequence of blocks defines integer prefixes by \(N_{n+1}=729N_n+\nu(b_{n+1})\), starting from \(N_0=m\), where \(m\) is the integer part determined by the readout. The preceding brackets place the closure, propagation, and self-scaling characters in \((3,4)\), \((2,3)\), and \((1,2)\), respectively; their initial values are therefore \(m=3,2,1\). Subsequent blocks refine the fractional part. The corresponding cylinder is
\begin{equation}
 I_n=\left[\frac{N_n}{729^n},\frac{N_n+1}{729^n}\right],
 \qquad |I_n|=729^{-n}.
 \label{eq:cilindro}
\end{equation}
The half-open convention is used to select a representative when a value lies on a boundary. Passage to the limit uses the closures of these intervals and identifies the two equivalent boundary expansions. This qualification prevents the assumption that every intersection of nested half-open intervals is automatically nonempty.

\begin{theorem}[Determination of the scalar realization]
A compatible block history determines a unique limit of the left endpoints of \eqref{eq:cilindro}. If an exact character has rational brackets of arbitrarily small width and is separated from the boundary of every decimal cylinder examined, each decimal prefix of finite length is determined after a finite computation. The prefixes thus obtained are mutually compatible.
\end{theorem}
\begin{proof}
Extending a block preserves the previous prefix by Euclidean division. The closed cylinders are nested and their diameter tends to zero; their left endpoints form a Cauchy sequence and have a unique limit. For a fixed decimal length, a value separated from the boundaries of the partition has positive distance from them. A sufficiently narrow bracket lies within a single cell and determines its prefix. Uniqueness of the cell and nesting of the partitions prove compatibility under truncation. At points with a terminating expansion, the exact boundary decision and its representation convention must be added; this decision does not follow from numerical width alone.
\end{proof}

\begin{corollary}[Finite determination and compatibility of regional prefixes]
\label{gen:prefijos-regionales}
For each of the realizations \(\pi,\varphi,e\) and every finite
decimal length, the constructed brackets determine the corresponding prefix
by a finite computation. The prefixes obtained at different depths
are compatible under truncation.
\end{corollary}
\begin{proof}
The closure, propagation, and self-scaling characters have the
preceding brackets. After their recognition, the irrationality of \(\pi,e,\varphi\)
ensures their separation from the rational boundaries of every finite decimal
partition. The preceding theorem applies to each character.
\end{proof}
Arbitrary depth expresses the quantifier in the corollary: for every
finite horizon there is a computation determining its prefix. The mathematical
object is the complete compatible collection; each execution obtains a
finite projection of it.

\subsection{Specular involution and synchronization of the regions}

The joint nonadic connection transports the three characters together. In the ninth-phase chart, the self-scaling axis and the ternary sum of closure and propagation remain fixed under specular interchange. Helical inversion changes the direction of transport; the mirror of the two channels changes their relative orientation. These are different involutions, although both act on the same trajectory system.

The matching census values illustrate this coordination:
\begin{center}
\begin{tabular}{@{}c rrr l@{}}\toprule
Phase&\(N_\pi\)&\(N_e\)&\(N_\varphi\)&Equality of counts\\\midrule
4&207&207&122&\(N_\pi=N_e\)\\
5&39&151&151&\(N_e=N_\varphi\)\\
6&110&110&69&\(N_\pi=N_e\)\\
7&81&29&81&\(N_\varphi=N_\pi\)\\
8&59&59&59&triple synchronization\\
9&43&43&43&triple synchronization\\\bottomrule
\end{tabular}
\end{center}
These integers count compatible coordinates of the regions. They do not assert equality between \(\pi\), \(e\), and \(\varphi\) as real numbers. If
\[
 \partial(a,b,c)=(a-b,b-c,c-a),
\]
its image belongs to the lattice \(A_2=\{(u,v,w)\in\ZZ^3:u+v+w=0\}\). In phases eight and nine the relative component vanishes, but the diagonal component changes from \((59,59,59)\) to \((43,43,43)\). Relative synchronization coexists with a shift of \(-16\) per channel. The closure that will lead to \(\alpha\) must preserve both pieces of information: relative coincidence and memory evolution.

